% =========================================================
% PARTE 6: CONCLUSIONES Y RESULTADOS OPERATIVOS
% =========================================================
\section{Parte 6: Conclusiones y Resultados Esperados}

\subsection{Respuestas a Preguntas Clave de Operación}

\begin{summarybox}{Respuestas a Preguntas Fundamentales de Negocio y Operación}
	\begin{enumerate}
		\item \textbf{¿Conviene seguir esperando a que la combi se llene con 18?}
		\textbf{No en la ruta Puente de Tuzos.} Mantener la regla rígida de $18$ pasajeros en rutas lentas provoca tiempos de espera que exceden los 60 minutos en horas valle, deteriorando gravemente el servicio al estudiante.
		\item \textbf{¿Cuál es el número mínimo razonable de pasajeros para salir?}
		En horas valle, salir con un piso de \textbf{12 a 14 pasajeros} tras 15--25 minutos de espera reduce el tiempo total acumulado de espera en más de un \textbf{40\%}, manteniendo una ocupación aceptable ($\ge 66\%$).
		\item \textbf{¿En qué horas conviene ser más flexible?}
		Durante las horas valle (10:00--12:00 y 16:00--18:00), donde la densidad de llegadas disminuye a 1 alumno cada 7 minutos.
		\item \textbf{¿Cuánto tiempo de espera se reduce saliendo con 12 o 14 personas?}
		Se ahorran en promedio \textbf{25 a 35 minutos de espera acumulada} por alumno en comparación con la política de combi llena.
	\end{enumerate}
\end{summarybox}

\subsection{Conclusión General}

El algoritmo UCB1 demuestra ser una herramienta efectiva para la optimización adaptativa del transporte universitario, logrando un balance equilibrado entre la satisfacción del estudiante (menor tiempo de espera) y la rentabilidad del concesionario de transporte (ocupación eficiente del vehículo).
